\documentclass[11pt, a4paper, oneside]{article}
\usepackage[ngerman]{babel}
\usepackage{worksheet}

\begin{document}
	\author{L. Bung}
	\title{Exkurs: Polynomdivision}
	\subject{Mathematik}
	\class{FHR1}
	\maketitle
	
	\grouptask*{Wiederholung: Schriftliche Division}
	
	a) Recherchieren Sie gemeinsam, wie die schriftliche Division nochmal ging.
	Beschreiben Sie das Verfahren anhand eines Beispiels im folgenden Kasten.
	
	\begin{warning}{Schriftliche Division}
		\phantom{x}\vspace{7cm}
	\end{warning}
	
	b) Nutzen Sie Ihre Erkenntnisse nun, um die folgenden Divisionsaufgaben zu lösen.
	
	\begin{checkeredfigure}
		\node at (1.25,6.75) {(i)};
		\node at (2.25,6.75) {84 : 4 =};
		\draw[thick] (1.5,6.5) -- (5,6.5);
		
		\node at (6.25,6.75) {(ii)};
		\node at (7.6,6.75) {512 : 8 =};
		\draw[thick] (6.5,6.5) -- (10.5,6.5);
		
		\node at (11.75,6.75) {(iii)};
		\node at (13.1,6.75) {503 : 9 =};
		\draw[thick] (12,6.5) -- (16,6.5);
	\end{checkeredfigure}
	
	Wir haben bereits gesehen, dass es Polynome $f(x) = a_nx^n+a_{n-1}x^{n-1}+\dots+a_1x+a_0$ gibt, in denen man Nullstellen nicht mit unseren bekannten Verfahren berechnen kann.
	Wenn wir allerdings schon eine Nullstelle $x_1$ wissen, können wir diese ja als Linearfaktor $(x-x_1)$ schreiben.
	
	Das Polynom hätte dann also die Form $f(x) = (x-x_1) \cdot r(x)$, wobei $r(x)$ ein Polynom vom Grad $n-1$ ist.
	Die Frage ist: Wie sieht das Polynom $r(x) = f(x) : (x-x_1)$ aus?
	
	\begin{hint}{Polynomdivision}
		Ähnlich wie bei der schriftlichen Division können wir das Restpolynom berechnen.
		Dieses Vorgehen nennt man \textbf{Polynomdivision}.
		Man geht dabei folgendermaßen vor:
		
		\begin{enumerate}
			\item Teile die führenden Terme und schreibe das Ergebnis hinter das Gleichheitszeichen.
			\item Multipliziere zurück und schreibe das Ergebnis unter das ursprüngliche Polynom.
			\item Subtrahiere den errechneten Term vom ursprünglichen Polynom.
			\item Wiederhole das Vorgehen, bis kein Rest mehr übrig bleibt.
		\end{enumerate}
	\end{hint}
	
	\task{Linearfaktoren durch Polynomdivision abspalten}
	
	Berechnen Sie jeweils die Restpolynome mithilfe der Polynomdivision.
	
	\begin{multicols}{2}
		\begin{enumerate}[label=\alph*)]
			\item $(x^2 - x - 6) : (x - 3)$
			\item $(x^3 + 2x^2 - x - 2) : (x + 1)$
		\end{enumerate}
	\end{multicols}
	
	\checkered[11cm]
	
	\task{Nullstellen mithilfe von Polynomdivision bestimmen}
	
	Wenn wir durch einen Linearfaktor teilen, hat das Restpolynom den Grad $n-1$.
	Das kann uns dabei helfen, weitere Nullstellen zu berechnen!
	
	Nutzen Sie die Polynomdivision, um alle Nullstellen zu bestimmen:
	
	\begin{multicols}{2}
		\begin{enumerate}[label=\alph*)]
			\item $(x^3 - 4x^2 + x + 6) : (x - 2)$
			\item $(2x^3 + 3x^2 - 11x - 6) : (x + 3)$
		\end{enumerate}
	\end{multicols}
	
	\checkered[13cm]
	
	\bonustask{Bedeutung des Rests}
	
	Führen Sie die Polynomdivision von $(x^3 + x + 1) : (x + 1)$ durch.
	
	\begin{itemize}
		\item Was fällt Ihnen in Bezug auf den Rest bei der Division auf?
		\item Welche Bedeutung könnte das haben?
	\end{itemize}
\end{document}